% !TEX program = xelatex
% ============================================================
%  Arabic Stock Option Plan Template — Nibras Center
%  خطة خيارات الأسهم للموظفين
%  Compile with: xelatex main.tex (twice)
% ============================================================

\documentclass[11pt,a4paper]{article}

\usepackage[a4paper,margin=2.5cm,top=2.5cm,bottom=2.5cm,headheight=15pt]{geometry}
\usepackage{array}
\usepackage{booktabs}
\usepackage{tabularx}
\usepackage{enumitem}
\usepackage{float}
\usepackage{titlesec}
\usepackage{fancyhdr}
\usepackage{setspace}
\usepackage[table]{xcolor}

\usepackage{bidi}
\usepackage[no-math]{fontspec}
\usepackage[quiet]{polyglossia}

\setmainfont[Scale=1]{Amiri-Regular.ttf}
\newfontfamily\arabicfont[Script=Arabic,Scale=0.95]{Amiri-Regular.ttf}
\newfontfamily\englishfont[Scale=0.95]{TeX Gyre Termes}

\setdefaultlanguage[calendar=gregorian,hijricorrection=1]{arabic}
\setotherlanguage[variant=british]{english}

\usepackage[breaklinks,hidelinks]{hyperref}
\hypersetup{pdftitle={خطة خيارات الأسهم},pdfauthor={الشركة}}

\titleformat{\section}{\large\bfseries}{\thesection}{0.7em}{}
\titleformat{\subsection}{\normalsize\bfseries}{\thesubsection}{0.6em}{}
\titlespacing*{\section}{0pt}{1.2ex plus 0.3ex}{0.8ex plus 0.2ex}

\pagestyle{fancy}
\fancyhf{}
\fancyhead[R]{\small\itshape خطة خيارات الأسهم}
\fancyhead[L]{\small اسم الشركة}
\fancyfoot[C]{\small\thepage}
\renewcommand{\headrulewidth}{0.3pt}

\newcolumntype{L}[1]{>{\raggedright\arraybackslash}p{#1}}
\newcolumntype{C}[1]{>{\centering\arraybackslash}p{#1}}

\definecolor{headerColor}{HTML}{1A3C6B}

\setstretch{1.4}

\begin{document}

\begin{center}
{\LARGE\bfseries خطة خيارات الأسهم للموظفين}\\[0.5em]
{\large (\LR{Employee Stock Option Plan — ESOP})}\\[1em]
{\normalsize اسم الشركة: \rule{6cm}{0.4pt}}\\[0.3em]
{\normalsize تاريخ الاعتماد: \rule{4cm}{0.4pt}}
\end{center}

\vspace{1.5em}

% ============================================================
% 1. PURPOSE
% ============================================================
\section{الغرض من الخطة}
\label{sec:purpose}
تهدف هذه الخطة إلى جذب الاحتفاظ بالكفاءات وتحفيز الموظفين على المساهمة في نمو الشركة من خلال منحهم خيارات شراء أسهم في الشركة.

% ============================================================
% 2. DEFINITIONS
% ============================================================
\section{التعريفات}
\label{sec:definitions}
\begin{itemize}
    \item \textbf{«الخطة»:} خطة خيارات الأسهم هذه.
    \item \textbf{«الشركة»:} \rule{6cm}{0.4pt}.
    \item \textbf{«الموظف»:} أي شخص يعمل لدى الشركة بدوام كامل.
    \item \textbf{«الخيار»:} حق شراء سهم واحد بسعر محدد.
    \item \textbf{«سعر التنفيذ»:} السعر المحدد لشراء السهم.
    \item \textbf{«الاستحقاق»:} اكتساب الحق في تنفيذ الخيار.
    \item \textbf{«التنفيذ»:} شراء الأسهم بسعر التنفيذ.
\end{itemize}

% ============================================================
% 3. POOL SIZE
% ============================================================
\section{حجم المجمع}
\label{sec:pool}
\begin{enumerate}
    \item يُخصّص مجمع خيارات إجمالي قدره \rule{3cm}{0.4pt} سهم، يمثل \rule{2cm}{0.4pt}\% من إجمالي أسهم الشركة.
    \item يجوز للجمعية العامة زيادة حجم المجمع بقرار مكتوب.
    \item تُدار الخطة من قبل مجلس الإدارة.
\end{enumerate}

% ============================================================
% 4. ELIGIBILITY
% ============================================================
\section{الأهلية}
\label{sec:eligibility}
\begin{enumerate}
    \item يحق للموظفين بدوام كامل الحصول على خيارات.
    \item لا يشمل ذلك الموظفين بدوام جزئي أو المقاولين المستقلين (إلا بقرار خاص).
    \item يُحدد مجلس الإدارة من يحصل على خيارات وكميتها.
\end{enumerate}

% ============================================================
% 5. GRANT
% ============================================================
\section{منح الخيارات}
\label{sec:grant}
\begin{enumerate}
    \item يُمنح الخيار بموجب «اتفاقية منح خيارات» (\LR{Grant Agreement}) يوقّعها الموظف والشركة.
    \item تحدد الاتفاقية: عدد الخيارات، سعر التنفيذ، تاريخ المنح، جدول الاستحقاق.
    \item يُحدد سعر التنفيذ بالقيمة العادلة للسهم في تاريخ المنح، كما يقدّره مجلس الإدارة.
\end{enumerate}

% ============================================================
% 6. VESTING
% ============================================================
\section{جدول الاستحقاق}
\label{sec:vesting}

\subsection{الجدول القياسي}
\begin{enumerate}
    \item \textbf{مدة الاستحقاق:} 4 سنوات من تاريخ المنح.
    \item \textbf{فترة الجرف (\LR{Cliff}):} سنة واحدة. لا تُستحق أي خيارات خلال السنة الأولى.
    \item \textbf{عند إكمال السنة الأولى:} تُستحق 25\% من الخيارات دفعة واحدة.
    \item \textbf{بعد الجرف:} تُستحق الخيارات المتبقية شهرياً (1/36 من الإجمالي كل شهر) على مدى 36 شهراً.
\end{enumerate}

\subsection{الاستحقاق المُسرّع}
\begin{enumerate}
    \item \textbf{استحواذ (\LR{Single-trigger}):} عند الاستحواذ على الشركة، تُستحق 50\% من الخيارات غير المستحقة فوراً.
    \item \textbf{استحواذ + إنهاء (\LR{Double-trigger}):} عند الاستحواذ وإنهاء الموظف خلال 12 شهراً، تُستحق جميع الخيارات فوراً.
\end{enumerate}

% ============================================================
% 7. EXERCISE
% ============================================================
\section{تنفيذ الخيارات}
\label{sec:exercise}
\begin{enumerate}
    \item يجوز للموظف تنفيذ الخيارات المستحقة في أي وقت.
    \item يُرسل إشعار تنفيذ مكتوب + دفع سعر التنفيذ.
    \item تُصدر الشركة الأسهم للموظف خلال \rule{2cm}{0.4pt} يوماً.
    \item \textbf{فترة التنفيذ بعد انتهاء العمل:}
    \begin{itemize}
        \item \textbf{استقالة طوعية:} \rule{2cm}{0.4pt} يوماً لتنفيذ الخيارات المستحقة.
        \item \textbf{إنهاء لسبب:} تُلغى جميع الخيارات (المستحقة وغير المستحقة).
        \item \textbf{وفاة أو عجز:} \rule{2cm}{0.4pt} شهراً للورثة لتنفيذ الخيارات المستحقة.
    \end{itemize}
\end{enumerate}

% ============================================================
% 8. TRANSFER RESTRICTIONS
% ============================================================
\section{قيود النقل}
\label{sec:transfer}
\begin{enumerate}
    \item لا يجوز نقل الخيارات أو الأسهم الناتجة عنها إلا بموافقة مجلس الإدارة.
    \item \textbf{حق الأولوية بالشراء (\LR{Right of First Refusal}):} للشركة الحق في شراء الأسهم قبل أي بيع لطرف ثالث.
    \item \textbf{قيد الاكتتاب العام (\LR{Lock-up}):} لا يجوز بيع الأسهم خلال \rule{2cm}{0.4pt} يوماً بعد الاكتتاب العام.
\end{enumerate}

% ============================================================
% 9. TAX
% ============================================================
\section{الأحكام الضريبية}
\label{sec:tax}
\begin{enumerate}
    \item يتحمل الموظف المسؤولية عن الضرائب الناتجة عن منح وتنفيذ الخيارات.
    \item قد تخضع الخيارات لمعاملة ضريبية خاصة وفقاً لقوانين \rule{4cm}{0.4pt}.
    \item يجوز للشركة حجز جزء من الأسومات لتغطية الضرائب.
\end{enumerate}

% ============================================================
% 10. AMENDMENT
% ============================================================
\section{تعديل الخطة}
\label{sec:amendment}
\begin{enumerate}
    \item يجوز لمجلس الإدارة تعديل الخطة، شريطة موافقة الجمعية العامة على التعديلات الجوهرية.
    \item لا يجوز تعديل شروط الخيارات الممنوحة بالفعل دون موافقة الموظف.
\end{enumerate}

% ============================================================
% 11. TERMINATION
% ============================================================
\section{إنهاء الخطة}
\label{sec:termination}
\begin{enumerate}
    \item يجوز لمجلس الإدارة إنهاء الخطة في أي وقت.
    \item لا يؤثر الإنهاء على الخيارات الممنوحة بالفعل.
    \item تُستحق جميع الخيارات عند الاكتتاب العام أو الاستحواذ.
\end{enumerate}

% ============================================================
% 12. SAMPLE GRANT
% ============================================================
\section{نموذج لاتفاقية منح خيارات}
\label{sec:sample}

\begin{table}[H]
\centering
\renewcommand{\arraystretch}{1.5}
\begin{tabularx}{\textwidth}{L{6cm} X}
\toprule
\textbf{البند} & \textbf{القيمة} \\
\midrule
اسم الموظف & --- \\
المنصب & --- \\
تاريخ المنح & --- \\
عدد الخيارات الممنوحة & --- \\
سعر التنفيذ للسهم & --- \LR{USD} \\
القيمة العادلة للسهم & --- \LR{USD} \\
جدول الاستحقاق & 4 سنوات، جرف سنة واحدة \\
فترة التنفيذ بعد الاستقالة & 90 يوماً \\
\bottomrule
\end{tabularx}
\end{table}

% ============================================================
% SIGNATURES
% ============================================================
\vspace{2em}
\section*{الاعتماد}

\begin{table}[H]
\centering
\renewcommand{\arraystretch}{2.5}
\begin{tabularx}{\textwidth}{L{5cm} X X}
\toprule
& \textbf{مجلس الإدارة} & \textbf{الجمعية العامة} \\
\midrule
\textbf{التاريخ} & \rule{4cm}{0.4pt} & \rule{4cm}{0.4pt} \\
\textbf{التوقيع} & \rule{4cm}{0.4pt} & \rule{4cm}{0.4pt} \\
\bottomrule
\end{tabularx}
\end{table}

\end{document}
